\section{Conclusiones}

\begin{enumerate}
    \item Los dos orígenes difieren más por su adquisición que por la clase de las
    imágenes. Su firma espectral está separada por completo y se concentra en la
    textura de escala fina (subbandas vertical y diagonal), no en la anatomía.
    \item Ninguna estrategia de realce mejoró la segmentación ni la imagen de forma
    consistente; \texttt{base\_adap}, sin realce, tuvo el mejor Dice en ambos
    orígenes, en entrenamiento y en validación.
    \item La segmentación desplaza los descriptores con tres tipos de error: sesgo
    sistemático, sesgo por sitio y error aleatorio. La relación de aspecto y la
    simetría son las menos reproducibles entre sitios.
    \item Con la máscara de referencia, \FinOECincoSigSH{} descriptores difieren entre
    grupos en Shenzhen; ninguno en Montgomery. Según la estrategia de segmentación
    cambian entre \FinDiscMin{} y \FinDiscMax{} de las 10 conclusiones: el
    preprocesamiento puede alterar las conclusiones científicas.
    \item Los descriptores automáticos estiman la fracción pulmonar de referencia en
    pacientes nuevos con un RMSE de \FinRMSE{} (IC 95\,\% \FinRMSEInf{} a
    \FinRMSESup{}), un \FinReduccion\,\% menos que la media de entrenamiento.
\end{enumerate}

\subsection{Trabajo futuro}

Comparar el error con la variabilidad entre observadores, para saber si es
clínicamente relevante; verificar los resultados en un tercer centro; usar máscaras
de referencia con una misma convención de trazado; y explorar otros parámetros de
realce y de filtrado.

\nocite{*}
\bibliographystyle{plain}
\bibliography{referencias}